\chapter{Local central charge and global residue}
\label{v4-arith:polyakov}

A local central charge is a coefficient in a stress-tensor
operator product. A pole order of a global scalar determinant
is a different invariant. Both can equal two in a rank-two
Heisenberg example, but their equality does not construct a
global Virasoro action.

\section{Rescaling and the conformal vector}
\label{v4-arith:pol:rescaling-not-conformal}

\begin{observation}[Rescaling and Virasoro normalization]
\label{v4-arith:obs:rescaling-breaks-virasoro}
Let \(T\) be a nonzero stress tensor on the Heisenberg algebra.
For \(T'=\lambda T\), its singular product is
\[
 T'(z)T'(w)\sim\lambda^2\left(
       \frac{c/2}{(z-w)^4}+\frac{2T(w)}{(z-w)^2}
                              +\frac{\partial T(w)}{z-w}\right).
\]
The standard conformal identities for the original translation
operator force \(\lambda=1\).
\end{observation}

The second pole requires \(\lambda^2=\lambda\).
If \(\lambda=0\), the zero state cannot implement the nonzero
translation operator. Thus the fourth-pole coefficient
\(\lambda^2c\) cannot be interpreted as a new central charge
while the other conformal identities are left unsatisfied.

\begin{remark}
At a prime \(p\), take the standard rank-two Sugawara state.
Its oscillator grading is \(L_0^{(p)}\), while the chosen
arithmetic energy is
\[
                         H_p=(\log p)L_0^{(p)}.
\]
The real number \(\log p\) weights this operator.
It does not rescale the conformal state.
\end{remark}

\begin{theorem}[The local central charge]
\label{v4-arith:thm:per-prime-c}
Each rank-two level-one Heisenberg factor has central charge two.
The operator \(H_p\) has energy \(m\log p\) on a
mode-\(m\) creator. It is not the zero mode of a scalar-rescaled
stress tensor satisfying the same conformal identities.
\end{theorem}

\begin{proof}
For two orthogonal currents, the stress tensor is
\(T=\tfrac12(:a^2:+:b^2:)\).
Its two independent double contractions give coefficient one
at the fourth pole, so \(c/2=1\).
The single contractions give the required second and first poles,
as in Proposition~\ref{v4-prop:singular-conformal-boundary}.
The energy formula follows from the oscillator grading.
The rescaling observation proves the final assertion.
\end{proof}

\section{The sum over primes}
\label{v4-arith:pol:infinite-tensor-not-VOA}

Write \(\mathcal H_{\mathbb P}=\bigotimes_p'\mathcal H^{(p)}\)
for the algebraic restricted tensor using the vacua.
A vector has nonvacuum factors at only finitely many primes.
This is the explicit current cover of
Chapter~\ref{v4-ch:singular-backgrounds}.
It is separate from any proposed arithmetic trace class.

\begin{observation}\label{v4-arith:obs:infinite-tensor-fails}
The formal sum \(\sum_p\omega_p\) is not a vector of
\(\mathcal H_{\mathbb P}\). In the Hilbert completion its
finite partial sums have squared norm equal to the number
of primes included, so they do not converge there either.
\end{observation}

The norm assertion follows because each \(\omega_p\) has
squared norm one and different prime sectors are orthogonal.
It rules out the indicated sum, with its indicated normalization.
The stronger absence of a conformal state implementing the
whole Hilbert oscillator grading is proved in
Proposition~\ref{v4-prop:singular-conformal-boundary}.

\begin{corollary}\label{v4-arith:cor:no-virasoro}
The sums of local modes \(\sum_pL_m^{(p)}\), for \(m\geq-1\),
act on \(\mathcal H_{\mathbb P}\). For \(m\leq-2\), these
same sums do not define operators on its vacuum.
Consequently the prescribed sums do not give a full
Virasoro representation on this carrier.
\end{corollary}

\begin{proof}
The modes with \(m\geq-1\) annihilate each local vacuum.
Only the finitely many occupied prime factors contribute.
For \(m\leq-2\), each local vacuum receives a nonzero
creation state. The infinite sum is not in the algebraic
restricted tensor and has divergent Hilbert norm.
\end{proof}

\section{The prime zeta function and analytic continuation}
\label{v4-arith:pol:prime-zeta-obstruction}

The prime zeta function is
\[
                      P(s)=\sum_p p^{-s},\qquad\operatorname{Re}s>1.
\]
Taking logarithms of the Euler product and applying Möbius
inversion gives
\begin{equation}\label{v4-eq:prime-zeta-mobius}
 \log\zeta(s)=\sum_{r\geq1}\frac{P(rs)}r,
 \qquad
 P(s)=\sum_{r\geq1}\frac{\mu(r)}r\log\zeta(rs).
\end{equation}
Here \(\mu\) is the Möbius function. The sums first converge
absolutely in the stated half-plane. Local branches of the
logarithms give the continuation where they are defined.

\begin{theorem}[Natural boundary of prime zeta]
\label{v4-arith:thm:prime-zeta-natural-boundary}
The line \(\operatorname{Re}s=0\) is a natural boundary
for the analytic continuation of \(P\).
In particular \(P\) has no holomorphic continuation to a
neighborhood of the origin.\footnote{C.-E.~Fr\"oberg,
\emph{On the prime zeta function}, BIT \textbf{8} (1968),
187--202. See also P.~L.~Krapivsky and J.~M.~Luck,
\emph{Zeta-regularization and natural boundaries},
\S2.2, equations~(2.11)--(2.12),
\href{https://arxiv.org/abs/2606.24536}{arXiv:2606.24536}.}
\end{theorem}

The real singularities at \(s=1/r\) occur when \(r\) is
squarefree. The coefficient \(\mu(r)\) matters:
\eqref{v4-eq:prime-zeta-mobius} does not assert a branch point
for every integer denominator or every putative zero contribution.
The natural-boundary theorem concerns the surviving singularities.

\begin{corollary}\label{v4-arith:cor:zeta-P-0-undefined}
Ordinary analytic continuation does not assign \(P(0)\)
a finite value. Hence it does not define a global central
charge by \(c=2P(0)\).
\end{corollary}

\begin{proof}
Such an evaluation would require a holomorphic germ at zero.
The natural-boundary theorem excludes that germ.
Other regularization prescriptions require their own definitions
and do not construct the missing stress-tensor modes.
\end{proof}

\begin{remark}
The value \(\zeta(0)=-1/2\) belongs to the zeta function
of all positive integers. Replacing \(P\) by \(\zeta\)
changes the indexing set of the regularized sum.
It does not evaluate the sum of the local prime charges.
\end{remark}

\section{Pole order and logarithmic residue}
\label{v4-arith:pol:c-as-residue}

A specified scalar determinant can have a finite pole order.
The meromorphic function itself, its leading Laurent coefficient,
and the residue of its logarithmic derivative must be distinguished.

\begin{theorem}\label{v4-arith:thm:regularised-c}
Suppose a scalar determinant is identified near \(s=1\) with
\(D(s)=\zeta(s)^2h(s)\), where \(h\) is holomorphic and
\(h(1)\ne0\). Then its pole order is two, and
\[
 -\operatorname{Res}_{s=1}\frac{D'(s)}{D(s)}=2,
 \qquad \lim_{s\to1}(s-1)^2D(s)=h(1).
\]
The separately defined scalar
\(r_\zeta=2\operatorname{Res}_{s=1}\zeta(s)\) also equals two.
None of these scalar calculations defines a Virasoro action.
\end{theorem}

\begin{proof}
Write \(\zeta(s)=(s-1)^{-1}u(s)\), where \(u\) is
holomorphic and \(u(1)=1\). Then
\[
 D(s)=(s-1)^{-2}u(s)^2h(s),\qquad
 \frac{D'}D=-\frac2{s-1}+2\frac{u'}u+\frac{h'}h.
\]
The last two terms are regular near one.
This proves the pole, limit and logarithmic-residue formulas.
The residue of \(\zeta\) is one, giving the stated value
of \(r_\zeta\).
\end{proof}

To apply the theorem to an arithmetic trace, one must construct
its determinant and prove the local factorization with \(h(1)\ne0\).
The equality \(r_\zeta=2\) alone proves neither premise.

\section{The surviving Witt modes}
\label{v4-arith:pol:witt-virasoro}

\begin{theorem}\label{v4-arith:thm:witt-virasoro}
For \(m,n\geq-1\) with \(m+n\geq-1\), the prescribed sums satisfy
\[
 [L_m,L_n]=(m-n)L_{m+n},
 \qquad L_m=\sum_pL_m^{(p)}.
\]
The bracket is zero when \(m=n=-1\).
The negative sums with index at most \(-2\) remain undefined
on the restricted vacuum.
\end{theorem}

\begin{proof}
On any vector, the contributing prime set is finite.
Different prime factors commute, so their Witt brackets add.
A Virasoro central term would require \(m+n=0\).
In this index range the possible values of \(m\) are
\(-1,0,1\), and \(m(m^2-1)=0\).
Thus the central term vanishes on these modes.
The vacuum calculation above excludes the remaining prescribed sums.
\end{proof}

\begin{remark}
For a finite prime set \(S\), the full local sum has central
charge \(2|S|\). When \(|m|\geq2\), its central commutator
coefficient grows with \(|S|\). There is no finite limit for
this unrenormalized coefficient. This is a statement about the
specified sums, not an obstruction to every other representation.
\end{remark}

\section{Mellin reflection and a modular comparison}
\label{v4-arith:pol:functional-equation-modular}

Set
\[
 \Lambda(s)=\pi^{-s/2}\Gamma(s/2)\zeta(s),\qquad
 \xi(s)=\tfrac12s(s-1)\Lambda(s).
\]
The first function retains the poles at zero and one.
The second removes them.

\begin{theorem}\label{v4-arith:thm:arithmetic-modular-invariance}
The completed zeta functions satisfy
\(\Lambda(s)=\Lambda(1-s)\) and \(\xi(s)=\xi(1-s)\).
If an independently specified scalar trace is identified with
one of these functions, it inherits that reflection with the
same normalization. Such an identification does not by itself
construct a genus-one state space or its modular action.
\end{theorem}

\begin{proof}
Let \(\theta(t)=\sum_{n\in\mathbb Z}e^{-\pi n^2t}\).
Poisson summation for the Gaussian gives
\(\theta(t)=t^{-1/2}\theta(1/t)\).
Its Mellin integral initially converges for \(\operatorname{Re}s>1\):
\[
 \Lambda(s)=\frac12\int_0^\infty(\theta(t)-1)t^{s/2}\frac{dt}{t}.
\]

Split at one and use the theta identity on \((0,1)\).
The result is
\[
 \Lambda(s)=-\frac1s-\frac1{1-s}
       +\frac12\int_1^\infty(\theta(t)-1)
                   \bigl(t^{s/2}+t^{(1-s)/2}\bigr)\frac{dt}{t}.
\]

The integral is entire by exponential decay and is invariant
under \(s\mapsto1-s\). The two rational terms exchange.
Multiplication by \(s(s-1)/2\) proves the assertion for \(\xi\).
Transport through a specified scalar identification gives the
conditional trace consequence.
\end{proof}

A vertex-algebra genus-one trace instead has the form
\(\operatorname{Tr}q^{L_0-c/24}\), on a specified graded state
space and convergence domain. Relating its modular substitution
to Mellin reflection requires a map between those constructions.
A shared functional equation does not supply that map.

\section{Arithmetic degree and a possible anomaly comparison}
\label{v4-arith:pol:arakelov-pairing}

For a finite divisor \(D=\sum_pn_p[p]\) on
\(\operatorname{Spec}\mathbb Z\), define its arithmetic degree by
\(\widehat{\deg}D=\sum_p n_p\log p\).
Archimedean metric contributions are separate data for a
metrized line bundle.

\begin{proposition}\label{v4-arith:prop:deligne-pairing}
The finite divisor \([p]\) has degree \(\log p\).
For \(\operatorname{Re}s>1\), these degree weights satisfy
\[
                  -\frac{\zeta'(s)}{\zeta(s)}
                       =\sum_p\sum_{r\geq1}(\log p)p^{-rs}.
\]
\end{proposition}

\begin{proof}
The first equality is the definition of arithmetic degree.
Differentiate the absolutely convergent logarithmic Euler product
\(\log\zeta(s)=\sum_{p,r}p^{-rs}/r\).
Normal convergence on smaller half-planes justifies the derivative.
\end{proof}

A conformal state \(\omega_p\) is a vector, not a metrized
line bundle. A Deligne pairing requires geometric line bundles.

For example, let \(\pi:Y\to S\) be a smooth projective curve
family of integral complex schemes, with line bundles \(L,M\).
Its Deligne pairing is the following determinant line.\footnote{
I.~Biswas, G.~Schumacher and L.~Weng,
\emph{Deligne pairing and determinant bundle}. Theorem~1,
\href{https://arxiv.org/abs/1106.0255}{arXiv:1106.0255}.}
\[
 \langle L,M\rangle=
 \det R\pi_*(L\otimes M)\otimes(\det R\pi_*L)^{-1}
 \otimes(\det R\pi_*M)^{-1}\otimes\det R\pi_*\mathcal O_Y.
\]

A metric refinement must additionally specify metrics and norms.
The scalar \(\log p\) does not identify this line with a
quadratic expression in the conformal state.

\begin{remark}
In the convention \(\langle T^\mu_\mu\rangle=-cR/(24\pi)\),
the two-dimensional Weyl anomaly describes metric variation of
a quantum field theory.\footnote{For the Euclidean effective action,
see T.~Bautista, A.~Benevides and A.~Dabholkar,
\emph{Nonlocal quantum effective actions in Weyl-flat spacetimes}.
Equations~(2.26) and~(3.2),
\href{https://arxiv.org/abs/1711.00135}{arXiv:1711.00135}.}
An arithmetic comparison would assign
metrized determinant lines to the relevant field data and compare
their variation with this functional. Arithmetic degree alone
does not prove such a variation formula.
\end{remark}

\section{The comparison data}
\label{v4-arith:pol:arakelov-zeta-structure}

The local current factors have central charge two.
Their restricted product has the prescribed Witt modes with
indices at least \(-1\), and its unchanged negative sums fail
on the vacuum. A scalar determinant with the specified
zeta factor has pole order two, with the logarithmic residue
computed above.

The arithmetic questions concern maps relating these objects.
A trace identification must preserve its Mellin normalization
and polar factors. A conformal or anomaly comparison must
additionally construct the state space, modes or metrized lines
whose operations are being compared.
